\clearpage
\section{Especificação formal M12: composição M12.1 $\times$ M12.2}
\label{sec:formal-m12}
\sloppy

M12.1 (Posts, comentários e históricos) e M12.2 (Roteiros, compartilhamento,
Storage e anexo a Post) foram provados isoladamente. Esta seção registra a
\textbf{prova composta}: um único modelo Alloy
(\texttt{specification/alloy/operations/composition\_m12.als}) com um
\texttt{EstadoIntegrado} onde o mesmo Post, vínculo canônico, professor, Roteiro,
anexo e comando representam os mesmos fatos ao longo da trajetória. Não se trata
de dois modelos justapostos: as transições são ligadas e os predicados de
autorização/coerência de fronteira são reproduzidos das fontes e verificados por
um guard de drift (\texttt{tools/formal/m12\_composition.mjs}). O contrato CUE já
fornece a ponte de dados aditiva (\texttt{\#M12\_2RoteiroAnexo} embute
\texttt{\#M12\_1RoteiroAnexo} e \texttt{\#M12\_2AnexoVinculado} liga o snapshot à
referência canônica); nenhum shape novo foi necessário.

\subsection*{Ponte de acesso e geração}
A ponte \texttt{ponteAcessoRoteiro} exige que todo acesso abstrato de M12.1
(\texttt{acessoRoteiroValidado}) seja fundamentado no acesso concreto de M12.2
(\texttt{acessoProfessorRoteiro}: dono ou compartilhamento atual, Roteiro
publicável com objeto e geração canônica). A publicação composta
(\texttt{publicarPostComAnexo}) exige ambos e prova
\texttt{x.aGeracao = geracaoValida[rot] = objetoGeracao[roteiroObj[rot]]}. Assim
\texttt{\#M12\_2RoteiroAnexo} refina a projeção M12.1 \textbf{sem} supor que
M12.1, isolado, já provava a geração
(\texttt{PublicacaoCompostaRefinaAcessoM12\_1},
\texttt{PublicacaoCompostaProvaGeracao}, \texttt{AnexoCompostoUsaGeracaoCanonica}).
Professor terceiro, Roteiro \texttt{PROVISORIO}/\texttt{VALIDADO} e geração
divergente não publicam (\texttt{TerceiroNaoPublicaComAnexo},
\texttt{ProvisorioNaoPublicaComAnexo}, \texttt{ValidadoNaoPublicaComAnexo},
\texttt{GeracaoDivergenteNaoPublicaComAnexo}).

\subsection*{Publicabilidade explícita (M12-CORR-01)}
O worklog corretivo de M12.2 afirmava que \texttt{compartilhar} chamava
\texttt{roteiroPublicavel[a,r]}, mas a proibição era obtida indiretamente por
\texttt{coerente[b]}. Nesta rodada a pré-condição tornou-se \textbf{explícita} em
\texttt{roteiros\_m12\_2.als}, alinhando o código à afirmação factual; o check
\texttt{CompartilharExigePublicavel} e o testemunho
\texttt{WitnessCompartilhaPublicavel} reexecutam a evidência.

\subsection*{Papel acadêmico versus vínculo persistido (M12-CORR-02)}
M12.2 excluía \texttt{vAluno in chefes} de \texttt{coerente} citando RN-ROLE-01,
mas RN-ROLE-01 proíbe acumular \emph{papéis}, enquanto
\texttt{Turma/\{id\}/Alunos/\{uid\}} é vínculo canônico com vida própria. M9/M11
\textbf{não} definem transição que elimine/vede atomicamente vínculos ao tornar-se
Chefe. A rota acadêmica passou a exigir o papel acadêmico autorizado
(\texttt{alunos}) \textbf{e} o vínculo canônico atual; a coerência não pressupõe
mais que um Chefe jamais reteve vínculo. Um Chefe com vínculo legado não usa a
rota acadêmica nem contorna Q13 (\texttt{RotaAcademicaExigePapelEVinculo},
\texttt{ChefeComVinculoLegadoNaoBaixaComposto} e a trajetória
\texttt{WitnessChefeComVinculoLegado}, que promove a Chefe preservando o vínculo e
prova a negação). A regra determinística Rust \texttt{aluno\_baixa} já exigia
\texttt{papel == "ALUNO"}; a mudança alinha Alloy, CUE (\texttt{papel}) e Rust.

\subsection*{Revogação Q09, aluno e Chefe}
A revogação preserva Post, snapshot, histórico, objeto e URLs já emitidas e impede
nova associação ou edição que mantém o anexo sem acesso
(\texttt{RevogacaoCompostaPreservaPost/Anexo/Objeto/Urls},
\texttt{RevogadoNaoMantemAnexoSemAcesso}); a desvinculação autorizada não exige
acesso e conserva o objeto (\texttt{DesvinculacaoNaoExigeAcesso},
\texttt{DesvinculoConservaObjeto}). O aluno com vínculo canônico atual baixa o
anexo de Post acessível, inclusive em turma arquivada
(\texttt{WitnessAlunoBaixaAtivo}, \texttt{WitnessAlunoBaixaArquivada}); ex-aluno
após remoção e refresh de claim não obtém nova URL
(\texttt{WitnessExAlunoNaoBaixa}); Post removido não aparece
(\texttt{WitnessPostRemovidoNaoBaixa}); URL já emitida só é utilizável enquanto
ativa (\texttt{WitnessUrlAtivaAposRevogacao}); claim antiga não autoriza
(\texttt{ClaimAntigaNaoAutorizaComposto}). O Chefe só emite sob escopo Q13
registrado, vinculado a Post removido que referencia o Roteiro
(\texttt{ChefeSemEscopoNaoEmiteComposto},
\texttt{WitnessChefeComEscopoQ13}); não cria Post alheio, não compartilha e não
usa a rota acadêmica (\texttt{ChefeNaoPublicaPostComposto},
\texttt{ChefeNaoCompartilhaComposto}, \texttt{ChefeAnexoSoQ13}); turma arquivada
nega escrita inclusive a intervenção (\texttt{TurmaArquivadaNegaEscritaComposta}).

\subsection*{M7/M9 na composição}
A primeira execução grava receipt sob o \texttt{idOperacao}
(\texttt{PrimeiraExecucaoCompostaProduzReceipt}); o retry idêntico não duplica
Post, anexo, histórico, compartilhamento nem notificação
(\texttt{RetryCompostoNaoDuplicaFato}); o reuso incompatível não herda receipt
(\texttt{ReusoIncompativelCompostoNaoHerda}); a revogação de vínculo impede o
commit posterior (\texttt{RevogacaoVinculoCompostaImpedeCommit}). Isto é
comportamento do modelo \emph{bounded}, não garantia de transação
Firestore--Storage nem revogação instantânea de URL já emitida.

\begingroup\small\raggedright
% Gerado por lcqui-spec-doc; não editar.
\subsection*{Evidência formal M12 --- composição M12.1 x M12.2 (Posts e Roteiros)}
Prova conjunta em um único EstadoIntegrado, com o mesmo Post, vínculo canônico, professor, Roteiro, anexo e comando. Reproduz os predicados de autorização/coerência de fronteira de M12.1/M12.2 (guard de drift em tools/formal/m12\_composition.mjs). A ponte `ponteAcessoRoteiro` exige que todo acesso abstrato M12.1 (`acessoRoteiroValidado`) seja fundamentado no acesso concreto M12.2 (`acessoProfessorRoteiro`, com objeto e geração canônica); a publicação composta prova a geração do anexo, ausente na prova isolada de M12.1. A rota acadêmica exige o papel acadêmico autorizado (`alunos`) E o vínculo canônico atual: um Chefe com vínculo legado (M9/M11 não removem vínculos atomicamente) não contorna Q13. Revogação Q09 preserva Post/snapshot/histórico; M7/M9 cobrem receipt, retry, reuso e revalidação. Escopos limitados a for 4/for 5. Não certifica Firebase, Storage real, Rules, auth, bytes nem atomicidade Firestore-Storage.
\begin{description}
\item[M12-INV-001] PublicacaoCompostaExigeDono: UNSAT.\newline Escopo: \texttt{check PublicacaoCompostaExigeDono for 4}.
\item[M12-INV-002] PublicacaoCompostaRefinaAcessoM12\_1: UNSAT.\newline Escopo: \texttt{check PublicacaoCompostaRefinaAcessoM12\_1 for 4}.
\item[M12-INV-003] SemAcessoM12\_1NaoPublicaComposto: UNSAT.\newline Escopo: \texttt{check SemAcessoM12\_1NaoPublicaComposto for 4}.
\item[M12-INV-004] PublicacaoCompostaProvaGeracao: UNSAT.\newline Escopo: \texttt{check PublicacaoCompostaProvaGeracao for 4}.
\item[M12-INV-005] AnexoCompostoUsaGeracaoCanonica: UNSAT.\newline Escopo: \texttt{check AnexoCompostoUsaGeracaoCanonica for 4}.
\item[M12-INV-006] TerceiroNaoPublicaComAnexo: UNSAT.\newline Escopo: \texttt{check TerceiroNaoPublicaComAnexo for 4}.
\item[M12-INV-007] ProvisorioNaoPublicaComAnexo: UNSAT.\newline Escopo: \texttt{check ProvisorioNaoPublicaComAnexo for 4}.
\item[M12-INV-008] ValidadoNaoPublicaComAnexo: UNSAT.\newline Escopo: \texttt{check ValidadoNaoPublicaComAnexo for 4}.
\item[M12-INV-009] GeracaoDivergenteNaoPublicaComAnexo: UNSAT.\newline Escopo: \texttt{check GeracaoDivergenteNaoPublicaComAnexo for 4}.
\item[M12-INV-010] CompartilhamentoCompostoSoPublicavel: UNSAT.\newline Escopo: \texttt{check CompartilhamentoCompostoSoPublicavel for 4}.
\item[M12-INV-011] AcessoCompartilhadoFundamentaM12\_1: UNSAT.\newline Escopo: \texttt{check AcessoCompartilhadoFundamentaM12\_1 for 4}.
\item[M12-INV-012] RevogacaoCompostaPreservaPost: UNSAT.\newline Escopo: \texttt{check RevogacaoCompostaPreservaPost for 4}.
\item[M12-INV-013] RevogacaoCompostaPreservaAnexo: UNSAT.\newline Escopo: \texttt{check RevogacaoCompostaPreservaAnexo for 4}.
\item[M12-INV-014] RevogacaoCompostaPreservaObjeto: UNSAT.\newline Escopo: \texttt{check RevogacaoCompostaPreservaObjeto for 4}.
\item[M12-INV-015] RevogacaoCompostaPreservaUrls: UNSAT.\newline Escopo: \texttt{check RevogacaoCompostaPreservaUrls for 4}.
\item[M12-INV-016] RevogadoNaoMantemAnexoSemAcesso: UNSAT.\newline Escopo: \texttt{check RevogadoNaoMantemAnexoSemAcesso for 4}.
\item[M12-INV-017] DesvinculacaoNaoExigeAcesso: UNSAT.\newline Escopo: \texttt{check DesvinculacaoNaoExigeAcesso for 4}.
\item[M12-INV-018] HistoricoAnexoCompostoPreservado: UNSAT.\newline Escopo: \texttt{check HistoricoAnexoCompostoPreservado for 4}.
\item[M12-INV-019] DesvinculoConservaObjeto: UNSAT.\newline Escopo: \texttt{check DesvinculoConservaObjeto for 4}.
\item[M12-INV-020] RotaAcademicaExigePapelEVinculo: UNSAT.\newline Escopo: \texttt{check RotaAcademicaExigePapelEVinculo for 4}.
\item[M12-INV-021] PostRemovidoNaoBaixaComposto: UNSAT.\newline Escopo: \texttt{check PostRemovidoNaoBaixaComposto for 4}.
\item[M12-INV-022] ChefeComVinculoLegadoNaoBaixaComposto: UNSAT.\newline Escopo: \texttt{check ChefeComVinculoLegadoNaoBaixaComposto for 4}.
\item[M12-INV-023] ClaimAntigaNaoAutorizaComposto: UNSAT.\newline Escopo: \texttt{check ClaimAntigaNaoAutorizaComposto for 4}.
\item[M12-INV-024] UrlEmitidaCompostaSobreviveRevogacao: UNSAT.\newline Escopo: \texttt{check UrlEmitidaCompostaSobreviveRevogacao for 4}.
\item[M12-INV-025] UrlExpiradaNaoUsavelComposto: UNSAT.\newline Escopo: \texttt{check UrlExpiradaNaoUsavelComposto for 4}.
\item[M12-INV-026] UrlAtivaUsavelComposto: UNSAT.\newline Escopo: \texttt{check UrlAtivaUsavelComposto for 4}.
\item[M12-INV-027] ChefeNaoPublicaPostComposto: UNSAT.\newline Escopo: \texttt{check ChefeNaoPublicaPostComposto for 4}.
\item[M12-INV-028] ChefeNaoCompartilhaComposto: UNSAT.\newline Escopo: \texttt{check ChefeNaoCompartilhaComposto for 4}.
\item[M12-INV-029] ChefeSemEscopoNaoEmiteComposto: UNSAT.\newline Escopo: \texttt{check ChefeSemEscopoNaoEmiteComposto for 4}.
\item[M12-INV-030] EscopoQ13CompostoExigePostRemovido: UNSAT.\newline Escopo: \texttt{check EscopoQ13CompostoExigePostRemovido for 4}.
\item[M12-INV-031] TurmaArquivadaNegaEscritaComposta: UNSAT.\newline Escopo: \texttt{check TurmaArquivadaNegaEscritaComposta for 4}.
\item[M12-INV-032] ChefeAnexoSoQ13: UNSAT.\newline Escopo: \texttt{check ChefeAnexoSoQ13 for 4}.
\item[M12-INV-033] PrimeiraExecucaoCompostaProduzReceipt: UNSAT.\newline Escopo: \texttt{check PrimeiraExecucaoCompostaProduzReceipt for 4}.
\item[M12-INV-034] RetryCompostoNaoDuplicaFato: UNSAT.\newline Escopo: \texttt{check RetryCompostoNaoDuplicaFato for 4}.
\item[M12-INV-035] ReusoIncompativelCompostoNaoHerda: UNSAT.\newline Escopo: \texttt{check ReusoIncompativelCompostoNaoHerda for 4}.
\item[M12-INV-036] RevogacaoVinculoCompostaImpedeCommit: UNSAT.\newline Escopo: \texttt{check RevogacaoVinculoCompostaImpedeCommit for 4}.
\item[M12-INV-037] PromoverChefeCompostoPreservaVinculo: UNSAT.\newline Escopo: \texttt{check PromoverChefeCompostoPreservaVinculo for 4}.
\item[M12-WIT-038] WitnessEstadoComposto: SAT.\newline Escopo: \texttt{run WitnessEstadoComposto for 4}.
\item[M12-WIT-039] WitnessTrajetoriaCompartilhaRevoga: SAT.\newline Escopo: \texttt{run WitnessTrajetoriaCompartilhaRevoga for 5}.
\item[M12-WIT-040] WitnessPublicaRoteiroCompartilhado: SAT.\newline Escopo: \texttt{run WitnessPublicaRoteiroCompartilhado for 4}.
\item[M12-WIT-041] WitnessProvisorioNaoPublica: SAT.\newline Escopo: \texttt{run WitnessProvisorioNaoPublica for 4}.
\item[M12-WIT-042] WitnessValidadoNaoPublica: SAT.\newline Escopo: \texttt{run WitnessValidadoNaoPublica for 4}.
\item[M12-WIT-043] WitnessAlunoBaixaAtivo: SAT.\newline Escopo: \texttt{run WitnessAlunoBaixaAtivo for 4}.
\item[M12-WIT-044] WitnessAlunoBaixaArquivada: SAT.\newline Escopo: \texttt{run WitnessAlunoBaixaArquivada for 4}.
\item[M12-WIT-045] WitnessExAlunoNaoBaixa: SAT.\newline Escopo: \texttt{run WitnessExAlunoNaoBaixa for 5}.
\item[M12-WIT-046] WitnessPostRemovidoNaoBaixa: SAT.\newline Escopo: \texttt{run WitnessPostRemovidoNaoBaixa for 4}.
\item[M12-WIT-047] WitnessRevogacaoPreservaHistorico: SAT.\newline Escopo: \texttt{run WitnessRevogacaoPreservaHistorico for 4}.
\item[M12-WIT-048] WitnessDesvinculaSemAcesso: SAT.\newline Escopo: \texttt{run WitnessDesvinculaSemAcesso for 4}.
\item[M12-WIT-049] WitnessChefeComEscopoQ13: SAT.\newline Escopo: \texttt{run WitnessChefeComEscopoQ13 for 5}.
\item[M12-WIT-050] WitnessChefeSemEscopoNaoEmite: SAT.\newline Escopo: \texttt{run WitnessChefeSemEscopoNaoEmite for 4}.
\item[M12-WIT-051] WitnessChefeComVinculoLegado: SAT.\newline Escopo: \texttt{run WitnessChefeComVinculoLegado for 4}.
\item[M12-WIT-052] WitnessUrlAtivaAposRevogacao: SAT.\newline Escopo: \texttt{run WitnessUrlAtivaAposRevogacao for 4}.
\item[M12-WIT-053] WitnessTurmaArquivadaNegaEscrita: SAT.\newline Escopo: \texttt{run WitnessTurmaArquivadaNegaEscrita for 4}.
\item[M12-WIT-054] WitnessRetryComposto: SAT.\newline Escopo: \texttt{run WitnessRetryComposto for 4}.
\item[M12-WIT-055] WitnessReusoIncompativelComposto: SAT.\newline Escopo: \texttt{run WitnessReusoIncompativelComposto for 4}.
\end{description}
UNSAT para check indica ausência de contraexemplo no escopo declarado; SAT para run indica testemunha encontrada. A evidência não certifica a implementação Firebase/backend.

\par\endgroup

\subsection*{Limites}
A prova composta usa escopos limitados (\texttt{for 4}/\texttt{for 5}) e não
certifica \texttt{functions/}, frontend, Rules, Firebase, Auth, Storage real,
bytes binários, HMAC, concorrência nem atomicidade entre serviços. Cada milestone
isolado (M12.1/M12.2) permanece válido e \textbf{não} é substituído por esta
seção; a propriedade composta é o que se demonstra aqui.
\fussy
